\begingroup
\renewcommand{\thepage}{B-\arabic{page}}
\setcounter{page}{1}
\pdfbookmark[0]{Individual appendix: hypothetical prompts}{individual-prompts}
\normalfont
\onehalfspacing
\raggedbottom
\setlength{\parindent}{0pt}
\setlength{\parskip}{0.25\baselineskip}

\section*{Appendix}

\begin{enumerate}
\setlength{\itemsep}{0.5\baselineskip}
\setlength{\parsep}{0pt}

\item I would like to investigate whether firms with more cash invest differently
in machinery and equipment after Hurricane Helene. How can I turn this into a
specific research question? Would it make sense to focus only on US manufacturing
firms?

\item I have chosen Duchin, Ozbas and Sensoy (2010) as my reference. Their study
focuses on the financial crisis, whereas mine concerns a hurricane. Can I apply
this research to a shock of this kind or use it to inform my study?

\item What data do I need to examine capex and cash in the context of Helene?
Should I use my own scraper for SEC filings in addition to Compustat?

\item You just suggested FEMA as a data source. What exactly is FEMA, what
information does it provide, and how can this help me assess firms' exposure
to Helene?

\item I would examine eight quarters before Helene and four full quarters
afterwards. Should I exclude 2024Q3 because Helene occurred at the end of
September?

\item I was planning to divide capex by current total assets. Could the
investment ratio then rise simply because assets are destroyed by the storm?
Would a fixed denominator from 2023 make sense? How would you implement this
in practice?

\item What would a difference-in-differences model that tests my question about
cash holdings look like? Please use Exposure, Post and Cash, along with firm
and quarter fixed effects. Explain each term and use only the methods covered
in our lecture slides.

\item Can I simplify the model without changing the research question?

\item I still do not understand the Post times Cash term. Why is the three-way
interaction alone not enough? Please explain what change among unexposed firms
we would otherwise miss.

\item Could firms with more cash also be better insured or better prepared?
To what extent do firm fixed effects account for these differences?

\item As an extension, I would subtract short-term debt from cash holdings.
How do I define this net liquidity measure using Compustat variables?
Should I retain negative values and use the same sample and model for the
comparison?

\item Please check my completed draft for inconsistencies between the research
question, hypothesis, variables and regression. Where can I cut repetition so
that the text fits within three pages? Suggest language improvements in English
without adding new results or further references.

\end{enumerate}
\clearpage
\endgroup
